%% Composition des vecteurs vitesse : un cascadeur (2) court sur le plateau
%% d'un camion (1) qui roule sur le sol (0).
%% V(M,2/0) = V(M,2/1) + V(M,1/0)
\documentclass[tikz,border=4pt]{standalone}
\usepackage{liaisons}
\begin{document}
\begin{tikzpicture}[x=1cm,y=1cm,
  vect/.style={-{Stealth[length=6pt]}, line width=1.5pt}]

% --- sol (0) ------------------------------------------------------------
\draw[black!55, line width=0.7pt] (-0.4,0) -- (6.6,0);
\node[font=\small, anchor=north east] at (6.55,-0.08) {$0$};

% --- camion plateau (1) -------------------------------------------------
\filldraw[fill=solideB!12, draw=solideB, line width=1.1pt, line join=round]
   (0.15,1.05) -- (0.15,2.30) -- (1.05,2.30) -- (1.28,1.80) -- (1.62,1.80) -- (1.62,1.05) -- cycle;
\filldraw[fill=solideB!12, draw=solideB, line width=1.1pt, line join=round]
   (1.62,1.05) -- (5.60,1.05) -- (5.60,0.72) -- (0.15,0.72) -- (0.15,1.05) -- cycle;
\draw[solideB, line width=0.9pt] (5.60,1.05) -- (5.60,1.55);   % ridelle arrière
\foreach \x in {1.05,4.55} {
  \filldraw[fill=black!70, draw=black!80, line width=0.8pt] (\x,0.42) circle (0.42);
  \filldraw[fill=black!25, draw=black!60, line width=0.6pt] (\x,0.42) circle (0.16);
}
\node[font=\small, text=solideB] at (0.72,1.55) {$1$};

% --- le cascadeur (2), qui court vers la droite -------------------------
%% le buste penché en avant dégage la verticale passant par M
\begin{scope}[line width=1.3pt, line cap=round, solideA]
  \draw[fill=white] (3.88,2.44) circle (0.19);            % tête
  \draw (3.55,1.66) -- (3.76,2.26);                       % tronc
  \draw (3.55,1.66) -- (3.24,1.06)                        % jambe arrière
        (3.55,1.66) -- (3.98,1.42) -- (4.16,1.06);        % jambe avant
  \draw (3.70,2.10) -- (3.36,1.94) (3.70,2.10) -- (4.12,2.26); % bras
\end{scope}
\fill[black] (3.55,1.66) circle (2pt);
\node[font=\small, anchor=east, inner sep=3pt] at (3.50,1.62) {$M$};
\node[font=\small, text=solideA] at (4.42,2.66) {$2$};

% --- les trois vecteurs vitesse au point M ------------------------------
%% tracés à trois hauteurs différentes pour rester lisibles ; le trait
%% pointillé rappelle qu'ils sont tous appliqués en M.
\draw[dotted, line width=0.7pt, black!60] (3.55,1.66) -- (3.55,4.30);
\draw[vect, solideA] (3.55,3.95) -- ++(-2.30,0)
      node[anchor=south, font=\small, text=solideA] at (2.40,4.00) {$\overrightarrow{V_{M,1/0}}$};
\draw[vect, solideB] (3.55,3.30) -- ++(1.45,0)
      node[anchor=south, font=\small, text=solideB] at (4.28,3.35) {$\overrightarrow{V_{M,2/1}}$};
\draw[vect, solideE] (3.55,2.68) -- ++(-0.85,0)
      node[anchor=south, font=\small, text=solideE] at (2.92,2.78) {$\overrightarrow{V_{M,2/0}}$};

% --- légende et relation -------------------------------------------------
\node[font=\small, anchor=north west, align=left] at (-0.4,4.35)
     {$0$ : sol\\ $1$ : camion\\ $2$ : cascadeur};
\node[draw=black!50, fill=yellow!25, rounded corners=2pt, line width=0.5pt,
      font=\small, anchor=north, inner sep=4pt] at (3.1,-0.30)
     {$\overrightarrow{V_{M,2/0}} = \overrightarrow{V_{M,2/1}} + \overrightarrow{V_{M,1/0}}$};
\end{tikzpicture}
\end{document}
